\documentclass[10pt,a4paper]{amsart}
\usepackage{amsmath,amssymb,amsthm,graphicx}
% Retained minimal preamble for the Stacks Project chapter sites-cohomology.tex (Cohomology on Sites).
% Verbatim excerpt of the chapter preamble: Xy-pic setup (native preamble lines 14-18) and the native
% theorem declarations with the mathematical macros (native preamble lines 154-217). Copied mechanically;
% nothing added, removed or redefined. The wrapper class and the amsmath/amssymb/amsthm/graphicx base
% packages are supplied by the assembler, so the class/externaldocument part of the native preamble is omitted.
% For commutative diagrams we use Xy-pic
\usepackage[all]{xy}

% We use 2cell for 2-commutative diagrams.
\xyoption{2cell}
% Theorem environments.
%
\theoremstyle{plain}
\newtheorem{theorem}[subsection]{Theorem}
\newtheorem{proposition}[subsection]{Proposition}
\newtheorem{lemma}[subsection]{Lemma}

\theoremstyle{definition}
\newtheorem{definition}[subsection]{Definition}
\newtheorem{example}[subsection]{Example}
\newtheorem{exercise}[subsection]{Exercise}
\newtheorem{situation}[subsection]{Situation}

\theoremstyle{remark}
\newtheorem{remark}[subsection]{Remark}
\newtheorem{remarks}[subsection]{Remarks}

\numberwithin{equation}{subsection}

% Macros
%
\def\lim{\mathop{\mathrm{lim}}\nolimits}
\def\colim{\mathop{\mathrm{colim}}\nolimits}
\def\Spec{\mathop{\mathrm{Spec}}}
\def\Hom{\mathop{\mathrm{Hom}}\nolimits}
\def\Ext{\mathop{\mathrm{Ext}}\nolimits}
\def\SheafHom{\mathop{\mathcal{H}\!\mathit{om}}\nolimits}
\def\SheafExt{\mathop{\mathcal{E}\!\mathit{xt}}\nolimits}
\def\Sch{\mathit{Sch}}
\def\Mor{\mathop{\mathrm{Mor}}\nolimits}
\def\Ob{\mathop{\mathrm{Ob}}\nolimits}
\def\Sh{\mathop{\mathit{Sh}}\nolimits}
\def\NL{\mathop{N\!L}\nolimits}
\def\CH{\mathop{\mathrm{CH}}\nolimits}
\def\proetale{{pro\text{-}\acute{e}tale}}
\def\etale{{\acute{e}tale}}
\def\QCoh{\mathit{QCoh}}
\def\Ker{\mathop{\mathrm{Ker}}}
\def\Im{\mathop{\mathrm{Im}}}
\def\Coker{\mathop{\mathrm{Coker}}}
\def\Coim{\mathop{\mathrm{Coim}}}

% Boxtimes
%
\DeclareMathSymbol{\boxtimes}{\mathbin}{AMSa}{"02}

%
% Macros for moduli stacks/spaces
%
\def\QCohstack{\mathcal{QC}\!\mathit{oh}}
\def\Cohstack{\mathcal{C}\!\mathit{oh}}
\def\Spacesstack{\mathcal{S}\!\mathit{paces}}
\def\Quotfunctor{\mathrm{Quot}}
\def\Hilbfunctor{\mathrm{Hilb}}
\def\Curvesstack{\mathcal{C}\!\mathit{urves}}
\def\Polarizedstack{\mathcal{P}\!\mathit{olarized}}
\def\Complexesstack{\mathcal{C}\!\mathit{omplexes}}
% \Pic is the operator that assigns to X its picard group, usage \Pic(X)
% \Picardstack_{X/B} denotes the Picard stack of X over B
% \Picardfunctor_{X/B} denotes the Picard functor of X over B
\def\Pic{\mathop{\mathrm{Pic}}\nolimits}
\def\Picardstack{\mathcal{P}\!\mathit{ic}}
\def\Picardfunctor{\mathrm{Pic}}
\def\Deformationcategory{\mathcal{D}\!\mathit{ef}}

% --- appended 2026-10-12 (Chief): the project own definitions of the reference/slogan/history
% display environments, copied mechanically from the shared Stacks preamble lines 8-12
% (\usepackage{verbatim} plus the three \newenvironment lines, which the chapter preamble keeps
% verbatim from the project preamble). Nothing added, removed or redefined: these environments
% carry no mathematical content and the project itself typesets them as comments. Without these
% three lines a counted statement whose own text uses them fails with Environment slogan undefined.
\usepackage{verbatim}
\newenvironment{reference}{\comment}{\endcomment}
\newenvironment{slogan}{\comment}{\endcomment}
\newenvironment{history}{\comment}{\endcomment}
\usepackage{cleveref}
\crefname{originalreference}{original reference}{original references}
\Crefname{originalreference}{Original reference}{Original references}
\makeatletter
\newlabel{sites-situation-inverse-limit-sites}{{0A34}{source}{}{sourcepointer-1}{}}
\newlabel{sites-situation-inverse-limit-sites@cref}{{[originalreference][0][]0A34}{[1][0][]source}{}{}{}}
\makeatother
\crefname{originallemma}{original lemma}{original lemmas}
\Crefname{originallemma}{Original lemma}{Original lemmas}
\makeatletter
\newlabel{sites-lemma-colimit-sites}{{09YL}{source}{}{sourcepointer-2}{}}
\newlabel{sites-lemma-colimit-sites@cref}{{[originallemma][0][]09YL}{[1][0][]source}{}{}{}}
\makeatother
\crefname{originallemma}{original lemma}{original lemmas}
\Crefname{originallemma}{Original lemma}{Original lemmas}
\makeatletter
\newlabel{sites-lemma-compute-pullback-to-limit}{{0A35}{source}{}{sourcepointer-3}{}}
\newlabel{sites-lemma-compute-pullback-to-limit@cref}{{[originallemma][0][]0A35}{[1][0][]source}{}{}{}}
\makeatother
\crefname{originallemma}{original lemma}{original lemmas}
\Crefname{originallemma}{Original lemma}{Original lemmas}
\makeatletter
\newlabel{sites-lemma-colimit}{{09YN}{source}{}{sourcepointer-4}{}}
\newlabel{sites-lemma-colimit@cref}{{[originallemma][0][]09YN}{[1][0][]source}{}{}{}}
\makeatother
\crefname{originallemma}{original lemma}{original lemmas}
\Crefname{originallemma}{Original lemma}{Original lemmas}
\makeatletter
\newlabel{lemma-colim-works-over-collection}{{0739}{source}{}{sourcepointer-5}{}}
\newlabel{lemma-colim-works-over-collection@cref}{{[originallemma][0][]0739}{[1][0][]source}{}{}{}}
\makeatother
\crefname{originallemma}{original lemma}{original lemmas}
\Crefname{originallemma}{Original lemma}{Original lemmas}
\makeatletter
\newlabel{lemma-cech-vanish-collection}{{03F9}{source}{}{sourcepointer-6}{}}
\newlabel{lemma-cech-vanish-collection@cref}{{[originallemma][0][]03F9}{[1][0][]source}{}{}{}}
\makeatother
\crefname{originallemma}{original lemma}{original lemmas}
\Crefname{originallemma}{Original lemma}{Original lemmas}
\makeatletter
\newlabel{lemma-injective-trivial-cech}{{03AW}{source}{}{sourcepointer-7}{}}
\newlabel{lemma-injective-trivial-cech@cref}{{[originallemma][0][]03AW}{[1][0][]source}{}{}{}}
\makeatother
\crefname{originallemma}{original lemma}{original lemmas}
\Crefname{originallemma}{Original lemma}{Original lemmas}
\makeatletter
\newlabel{algebra-lemma-directed-colimit-exact}{{00DB}{source}{}{sourcepointer-8}{}}
\newlabel{algebra-lemma-directed-colimit-exact@cref}{{[originallemma][0][]00DB}{[1][0][]source}{}{}{}}
\makeatother
\title{Cohomology of a colimit of sheaves over a colimit of sites as the colimit of the cohomologies}
\author{The Stacks Project Authors}
\date{stacks-project:sites-cohomology.tex:local-snapshot-20261011; GNU FDL 1.2 or later}
\begin{document}
\maketitle
\noindent\textit{Editorial note.} All delivered mathematical text is the author's own, copied from the pinned Stacks Project chapter file sites-cohomology.tex; no statement or proof step was rewritten, reordered or removed.\par\medskip
\noindent\textit{Editorial note.} The proof is the author's complete argument for all degrees; the reduction is modelled on the companion result of the same chapter and the Cech computation is carried out here.\par\medskip
\noindent\textit{Editorial note.} Every labelled reference occurring inside the counted spans is carried in assembly.external\_references as a non-load-bearing navigation pointer (same-chapter ones at their definition site, cross-chapter ones at the invocation inside the counted span), so no load-bearing mathematics is replaced by a stub. No bibliography citation occurs inside the counted spans.\par\medskip
\noindent\textit{Editorial note.} The checked chapter list for same-name collisions is the set of stacks chapters staged in this repository: chow, cohomology, divisors, etale-cohomology, flat, more-algebra, more-morphisms, morphisms, spaces-more-morphisms, varieties together with this card's chapter sites-cohomology.\par\medskip
\noindent\textit{Editorial note.} The preamble is the chapter's own retained minimal preamble (sources/stacks-sites-cohomology/retained-preamble.tex), a verbatim excerpt of the native preamble: the Xy-pic setup and the native theorem declarations with the mathematical macros.\par\medskip
\noindent\textit{Allowed standard prerequisites (editorial).}\begin{itemize}
\item The inverse system of sites and the sheaves attached to it. Same chapter, the statement of lemma-colim-sites-injective delivered verbatim as a context span in this card (source lines 2854-2875), fixing the cofinal index category, the sites, the sheaves and their compatibility maps. Supplies every object appearing in the counted statement and the embedding used in the proof.
\item Cech cohomology vanishes for the members of the system and filtered colimits of acyclic complexes are acyclic. Same chapter, the vanishing of Cech cohomology of injective sheaves, the criterion computing cohomology by Cech cohomology over a cofinal collection of coverings, and the exactness of filtered colimits. Supplies the three ingredients of the Cech computation.
\end{itemize}\medskip
\ifdefined\hypertarget\hypertarget{sourcepointer-1}{}\fi
\noindent\textit{Source reference.} Original reference 0A34: sources/stacks-sites-cohomology/pinned-original-source.tex, source lines 2858--2858. Navigation pointer invoked at proof line 2858; the counted span contains the pointer only and the cited result is not reproduced inside it, so the pointer is registered as a non-load-bearing navigation stub in the established form for Stacks material. No bibliography citation is involved.\par\medskip
\ifdefined\hypertarget\hypertarget{sourcepointer-2}{}\fi
\noindent\textit{Source reference.} Original lemma 09YL: sources/stacks-sites-cohomology/pinned-original-source.tex, source lines 2860--2860. Navigation pointer invoked at proof line 2860; the counted span contains the pointer only and the cited result is not reproduced inside it, so the pointer is registered as a non-load-bearing navigation stub in the established form for Stacks material. No bibliography citation is involved.\par\medskip
\ifdefined\hypertarget\hypertarget{sourcepointer-3}{}\fi
\noindent\textit{Source reference.} Original lemma 0A35: sources/stacks-sites-cohomology/pinned-original-source.tex, source lines 2861--2861. Navigation pointer invoked at proof line 2861; the counted span contains the pointer only and the cited result is not reproduced inside it, so the pointer is registered as a non-load-bearing navigation stub in the established form for Stacks material. No bibliography citation is involved.\par\medskip
\ifdefined\hypertarget\hypertarget{sourcepointer-4}{}\fi
\noindent\textit{Source reference.} Original lemma 09YN: sources/stacks-sites-cohomology/pinned-original-source.tex, source lines 2917--2917. Navigation pointer invoked at proof line 2917; the counted span contains the pointer only and the cited result is not reproduced inside it, so the pointer is registered as a non-load-bearing navigation stub in the established form for Stacks material. No bibliography citation is involved.\par\medskip
\ifdefined\hypertarget\hypertarget{sourcepointer-5}{}\fi
\noindent\textit{Source reference.} Original lemma 0739: sources/stacks-sites-cohomology/pinned-original-source.tex, source lines 2518--2518. Navigation pointer invoked at proof line 2923; the counted span contains the pointer only and the cited result is not reproduced inside it, so the pointer is registered as a non-load-bearing navigation stub in the established form for Stacks material. No bibliography citation is involved.\par\medskip
\ifdefined\hypertarget\hypertarget{sourcepointer-6}{}\fi
\noindent\textit{Source reference.} Original lemma 03F9: sources/stacks-sites-cohomology/pinned-original-source.tex, source lines 1565--1565. Navigation pointer invoked at proof line 2932; the counted span contains the pointer only and the cited result is not reproduced inside it, so the pointer is registered as a non-load-bearing navigation stub in the established form for Stacks material. No bibliography citation is involved.\par\medskip
\ifdefined\hypertarget\hypertarget{sourcepointer-7}{}\fi
\noindent\textit{Source reference.} Original lemma 03AW: sources/stacks-sites-cohomology/pinned-original-source.tex, source lines 1334--1334. Navigation pointer invoked at proof line 2942; the counted span contains the pointer only and the cited result is not reproduced inside it, so the pointer is registered as a non-load-bearing navigation stub in the established form for Stacks material. No bibliography citation is involved.\par\medskip
\ifdefined\hypertarget\hypertarget{sourcepointer-8}{}\fi
\noindent\textit{Source reference.} Original lemma 00DB: sources/stacks-sites-cohomology/pinned-original-source.tex, source lines 2943--2943. Navigation pointer invoked at proof line 2943; the counted span contains the pointer only and the cited result is not reproduced inside it, so the pointer is registered as a non-load-bearing navigation stub in the established form for Stacks material. No bibliography citation is involved.\par\medskip
\begin{lemma}
\label{lemma-colim-sites-injective}
Let $\mathcal{I}$ be a cofiltered index category and let
$(\mathcal{C}_i, f_a)$ be an inverse system of sites over $\mathcal{I}$
as in Sites, Situation \ref{sites-situation-inverse-limit-sites}.
Set $\mathcal{C} = \colim \mathcal{C}_i$ as in Sites,
Lemmas \ref{sites-lemma-colimit-sites} and
\ref{sites-lemma-compute-pullback-to-limit}.
Moreover, assume given
\begin{enumerate}
\item an abelian sheaf $\mathcal{F}_i$ on $\mathcal{C}_i$ for all
$i \in \Ob(\mathcal{I})$,
\item for $a : j \to i$ a map
$\varphi_a : f_a^{-1}\mathcal{F}_i \to \mathcal{F}_j$
of abelian sheaves on $\mathcal{C}_j$
\end{enumerate}
such that $\varphi_c = \varphi_b \circ f_b^{-1}\varphi_a$
whenever $c = a \circ b$. Then there exists a map of systems
$(\mathcal{F}_i, \varphi_a) \to (\mathcal{G}_i, \psi_a)$
such that $\mathcal{F}_i \to \mathcal{G}_i$ is injective and
$\mathcal{G}_i$ is an injective abelian sheaf.
\end{lemma}

\begin{lemma}
\label{lemma-colimit}
In the situation of Lemma \ref{lemma-colim-sites-injective} set
$\mathcal{F} = \colim f_i^{-1}\mathcal{F}_i$.
Let $i \in \Ob(\mathcal{I})$, $X_i \in \text{Ob}(\mathcal{C}_i)$. Then
$$
\colim_{a : j \to i} H^p(u_a(X_i), \mathcal{F}_j) =
H^p(u_i(X_i), \mathcal{F})
$$
for all $p \geq 0$.
\end{lemma}

\begin{proof}
The case $p = 0$ is Sites, Lemma \ref{sites-lemma-colimit}.

\medskip\noindent
Choose $(\mathcal{F}_i, \varphi_a) \to (\mathcal{G}_i, \psi_a)$
as in Lemma \ref{lemma-colim-sites-injective}.
Arguing exactly as in the proof of
Lemma \ref{lemma-colim-works-over-collection}
we see that it suffices to prove that
$H^p(X, \colim f_i^{-1}\mathcal{G}_i) = 0$ for $p > 0$.

\medskip\noindent
Set $\mathcal{G} = \colim f_i^{-1}\mathcal{G}_i$.
To show vanishing of cohomology of $\mathcal{G}$ on every object
of $\mathcal{C}$ we show that the {\v C}ech cohomology of $\mathcal{G}$
for any covering $\mathcal{U}$ of $\mathcal{C}$ is zero
(Lemma \ref{lemma-cech-vanish-collection}).
The covering $\mathcal{U}$ comes from a covering
$\mathcal{U}_i$ of $\mathcal{C}_i$ for some $i$. We have
$$
\check{\mathcal{C}}^\bullet(\mathcal{U}, \mathcal{G}) =
\colim_{a : j \to i}
\check{\mathcal{C}}^\bullet(u_a(\mathcal{U}_i), \mathcal{G}_j)
$$
by the case $p = 0$. The right hand side is acyclic in positive degrees
as a filtered colimit of acyclic complexes by
Lemma \ref{lemma-injective-trivial-cech}. See
Algebra, Lemma \ref{algebra-lemma-directed-colimit-exact}.
\end{proof}

\end{document}
